% Part:sets-functions-relations
% Chapter: size-of-sets
% Section: zig-zag

\documentclass[../../../include/open-logic-section]{subfiles}

\begin{document}

\olfileid{sfr}{siz}{zigzag}

\olsection{कँटर यांची नागमोडी पद्धत}

\begin{explain}
काही ``सोप्या'' प्रगणनांचा आपण आधीच विचार केला आहे. आता जरा कठीण
उदाहरण पाहू. नैसर्गिक संख्यांच्या जोड्यांचा संच विचारात घ्या
\oliflabeldef{sfr:set:pai:sec}{; त्याची व्याख्या आपण
\olref[set][pai]{sec} मध्ये पुढीलप्रमाणे केली:}{; त्याची व्याख्या पुढीलप्रमाणे आहे:}
\[
\Nat \times \Nat = \Setabs{\tuple{n,m}}{n,m \in \Nat}
\]
या क्रमित जोड्या आपण पुढीलप्रमाणे एका \emph{सरणीत} मांडू शकतो:
\[
\begin{array}{ c | c | c | c | c | c}
& \mathbf 0 & \mathbf 1 & \mathbf 2 & \mathbf 3 & \dots \\
\hline
\mathbf 0 & \tuple{0,0} & \tuple{0,1} & \tuple{0,2} & \tuple{0,3} & \dots \\
\hline
\mathbf 1 & \tuple{1,0} & \tuple{1,1} & \tuple{1,2} & \tuple{1,3} & \dots \\
\hline
\mathbf 2 & \tuple{2,0} & \tuple{2,1} & \tuple{2,2} & \tuple{2,3} & \dots \\
\hline
\mathbf 3 & \tuple{3,0} & \tuple{3,1} & \tuple{3,2} & \tuple{3,3} & \dots \\
\hline
\vdots & \vdots & \vdots & \vdots & \vdots & \ddots\\
\end{array}
\]
$\Nat \times \Nat$ मधील प्रत्येक क्रमित जोडी या सरणीत नेमकी एकदाच
येईल, हे स्पष्ट आहे. विशेषतः, $\tuple{n,m}$ ही जोडी $n$ व्या ओळीत
आणि $m$ व्या स्तंभात येईल. पण अशा सरणीतील घटकांची ``एकमितीय'' यादी
कशी करायची? पुढील सरणीतील नमुना ती करण्याचा एक मार्ग दाखवतो
(अर्थातच इतर अनेक मार्गही आहेत):
\[
\begin{array}{ c | c | c | c | c | c | c}
& \mathbf 0 & \mathbf 1 & \mathbf 2 & \mathbf 3 & \mathbf 4 &\dots \\
\hline
\mathbf 0 & 0  & 1& 3 & 6& 10 &\ldots \\
\hline
\mathbf 1 &2 & 4& 7 & 11 & \dots &\ldots \\
\hline
\mathbf 2 & 5 & 8 & 12 & \ldots & \dots&\ldots \\
\hline
\mathbf 3 & 9 & 13 & \ldots & \ldots & \dots & \ldots \\
\hline
\mathbf 4 & 14 & \ldots & \ldots & \ldots & \dots & \ldots \\
\hline
\vdots & \vdots & \vdots & \vdots & \vdots&\ldots & \ddots\\
\end{array}
\]\noindent
या नमुन्याला \emph{कँटर यांची नागमोडी पद्धत} म्हणतात. तो
$\Nat \times \Nat$ चे पुढीलप्रमाणे प्रगणन करतो:
\[
\tuple{0,0}, \tuple{0,1}, \tuple{1,0}, \tuple{0,2}, \tuple{1,1},
\tuple{2,0}, \tuple{0,3}, \tuple{1,2}, \tuple{2,1}, \tuple{3,0}, \dots
\]
यावरून पुढील विधान सिद्ध होते:
\end{explain}

\begin{prop}\ollabel{natsquaredenumerable}
$\Nat \times \Nat$ हा !!{enumerable} आहे.
\end{prop}

\begin{proof}
प्रत्येक $k \in \Nat$ ला कँटर यांच्या नागमोडी सरणीतील ज्या $n$ व्या
ओळीच्या आणि $m$ व्या स्तंभाच्या जागेचे मूल्य $k$ आहे, त्या
$\tuple{n,m} \in \Nat \times \Nat$ या क्रमित घटकाशी जोडणारे
$f \colon \Nat \to \Nat\times\Nat$ असे फलन घ्या.
\end{proof}

\begin{explain}
ही पद्धत अधिक व्यापक बाबतीतही सहज लागू होते. उदाहरणार्थ, नैसर्गिक
संख्यांच्या क्रमित तीन-घटकसमूहांचा पुढील संच प्रगणित करण्यासाठी ती
वापरता येते:
\[
\Nat \times \Nat \times \Nat = \Setabs{\tuple{n,m,k}}{n,m,k \in \Nat}
\]
$\Nat \times \Nat \times \Nat$ हा $\Nat \times \Nat$ आणि $\Nat$
यांचा कार्तीय गुणाकार आहे असे आपण मानतो; म्हणजे,
\[
\Nat^3 = (\Nat \times \Nat) \times \Nat =
\Setabs{\tuple{\tuple{n,m},k}}{n, m, k
  \in \Nat }
\]
म्हणून एका अक्षाला $\Nat$ चे प्रगणन आणि दुसऱ्या अक्षाला $\Nat^2$
चे प्रगणन अशी नावे देऊन, एका सरणीच्या साहाय्याने $\Nat^3$ चे
प्रगणन करता येते:
\[
\begin{array}{ c | c | c | c | c | c}
& \mathbf 0 & \mathbf 1 & \mathbf 2 & \mathbf 3 & \dots \\
\hline
\mathbf{\tuple{0,0}} & \tuple{0,0,0} & \tuple{0,0,1} & \tuple{0,0,2} & \tuple{0,0,3} & \dots \\
\hline
\mathbf{\tuple{0,1}} & \tuple{0,1,0} & \tuple{0,1,1} & \tuple{0,1,2} & \tuple{0,1,3} & \dots \\
\hline
\mathbf{\tuple{1,0}} & \tuple{1,0,0} & \tuple{1,0,1} & \tuple{1,0,2} & \tuple{1,0,3} & \dots \\
\hline
\mathbf{\tuple{0,2}} & \tuple{0,2,0} & \tuple{0,2,1} & \tuple{0,2,2} & \tuple{0,2,3} & \dots\\
\hline
\vdots & \vdots & \vdots & \vdots & \vdots & \ddots \\
\end{array}
\]
अशा रीतीने कँटर यांच्या नागमोडी पद्धतीसारखी पद्धत वापरून $\Nat^3$
चे प्रगणन मिळते. ही प्रक्रिया पुढे चालू ठेवून, प्रत्येक नैसर्गिक संख्या
$n$ साठी $\Nat^n$ चे प्रगणन मिळवता येते. म्हणून पुढील विधान मिळते:
\end{explain}

\begin{prop}
प्रत्येक $n \in \Nat$ साठी $\Nat^n$ हा !!{enumerable} आहे.
\end{prop}

\begin{prob}\label{sfr:siz:zigzag:prob:posint-n}
प्रत्येक $n \in \Nat$ साठी $(\PosInt)^n$ हा !!{enumerable} आहे, हे दाखवा.
\end{prob}

\begin{prob}\label{sfr:siz:zigzag:prob:posint-star}
  $(\PosInt)^*$ हा !!{enumerable} आहे, हे दाखवा. तुम्ही \cref{sfr:siz:zigzag:prob:posint-n} गृहीत धरू शकता.
\end{prob}

\end{document}
